\begin{tikzpicture}[figure picture,x=10mm,y=10mm,font=\figurefont\fontsize{7.6}{9.6}\selectfont]
\path[use as bounding box] (0,0) rectangle (16.4,4.9);
% 公共底板：三条横线区间
\foreach \xa/\xb in {0.5/5.1, 5.8/10.4, 11.1/15.9} \draw[FigureGrid,line width=.8pt] (\xa,1.4)--(\xb,1.4);
% (a) 绝对连续：μ 均匀铺底，ν 是密度曲线，无点质量
\node[figure heading,anchor=west] at (0.5,4.55) {(a) 绝对连续 $\nu\ll\mu$};
\fill[FigureBlueFill] (0.5,1.34) rectangle (5.1,1.46);
\fill[FigureRedFill] plot[domain=0.9:4.7,samples=60] (\x,{1.4+1.35*exp(-2*(\x-2.8)*(\x-2.8))}) |- (0.9,1.4) -- cycle;
\draw[FigureRed,line width=1pt,domain=0.9:4.7,samples=60] plot (\x,{1.4+1.35*exp(-2*(\x-2.8)*(\x-2.8))});
\node[left,FigureBlue] at (0.55,1.4) {$\mu$};
\node[right,FigureRed] at (4.75,2.1) {$\nu$};
\node[figure note,align=center] at (2.8,.55) {$\mu$ 为零的地方，$\nu$ 也为零；\\$\nu$ 可写成密度，没有点质量};
% (b) 等价：两条都处处为正
\node[figure heading,anchor=west] at (5.8,4.55) {(b) 等价 $\mu\sim\nu$};
\draw[FigureBlue,line width=1pt,domain=6.2:10.0,samples=60] plot (\x,{1.4+0.75+0.22*sin(120*\x)});
\draw[FigureRed,line width=1pt,domain=6.2:10.0,samples=60] plot (\x,{1.4+0.42+0.14*sin(120*\x+40)});
\node[left,FigureBlue] at (6.25,2.18) {$\mu$};
\node[left,FigureRed] at (6.25,1.8) {$\nu$};
\node[figure note,align=center] at (8.1,.55) {两个方向都绝对连续，\\零集完全相同（但质量不必相等）};
% (c) 相互奇异：连续长度与点质量分居
\node[figure heading,anchor=west] at (11.1,4.55) {(c) 相互奇异 $\mu\perp\nu$};
\fill[FigureBlueFill] (11.1,1.34) rectangle (13.35,1.46);
\fill[FigureBlueFill] (13.6,1.34) rectangle (15.9,1.46);
\fill[FigureRed] (13.47,1.4) circle[radius=1.1mm];
\node[left,FigureBlue] at (11.15,1.4) {$\mu$};
\node[above,FigureRed] at (13.47,1.62) {$\nu=\delta_0$};
\node[below,FigureMuted] at (13.47,1.18) {$0$};
\node[figure note,align=center] at (13.5,.55) {$\mu$ 的质量在补集上，$\nu$ 的质量在 $\{0\}$ 上，\\两部分可测意义下互不交叠};
\end{tikzpicture}
